%% infotel.tex 1.4   2016-09-15    Infotel latex template
%------------------------------------------------------------------
% Filename: infotel_template.tex
%
% This file is intended as a template for typesetting articles for the
%
%                        Journal of Spatial Information Science.
%
% Please edit this template to generate your own formatted manuscripts




%%% * All titles and sections lower case *EXCEPT short title  [ ]
%%% * Remove author postal addresses, only have geographic places and institutions [ ] 
%%% * Consistent use of Section, Figure, Table (capitalized and in full) [ ]
%%% * 10 keywords (and all lower case) [ ]
%%% * Remove all avoidable footnotes [ ]
%%% * Use double quotation marks (``'' not "" or `') [ ]
%%% * Punctuation inside quotations [ ]
%%% * E.g. and i.e. followed by comma [ ]
%%% * cf. followed by tilde [ ]
%%% * Itemize and enumerate correctly punctuated [e.g., "1. x, 2. y, and 3. x." ]
%%% * And/or lists using American English punctuation (e.g., "x, y, and z") [ ] 
%%% * Bibliography (e.g., en-dashes for number ranges, consistent "Proc.~" for Proceedings of..., etc.) []
%%% * Acknowledgment style use section* [ ] 
%%% * et al. no italics, but with dot  [ ] 
%%% * All captions end with full stop  [ ] 
%%% * Table captions under, not over table  [ ]
%%% * Adjust urls with burlalt [ ] 
%%% * Check correct use of hyphens, emdashes, endashes  [ ]
%%% * Perform spell check  [ ] 


%%% Check DOI, page numbers on article and web site. [ ]
%%% Update web site with final title, abstract, keywords. [ ] 
%%% Build with distiller for DOI links. [ ]


\documentclass{infotel}
\usepackage{hyperref}
\hypersetup{
	colorlinks=false,
	citecolor=Violet,
	linkcolor=Red,
	urlcolor=Blue}
\usepackage[hyphenbreaks]{breakurl}
\usepackage{booktabs}
\usepackage{stmaryrd}
\usepackage[T1]{fontenc}
\usepackage{cite}
\usepackage{hyperref}
\hypersetup{colorlinks,linkcolor={blue},citecolor={blue},urlcolor={red}}  

% Suggested packages for algorithm formatting
\usepackage{algorithm}
%\usepackage{algorithmic}
\usepackage{algpseudocode}


\usepackage[table]{xcolor}
\usepackage{amssymb,amsmath}

%\setcounter{page}{412}
\renewcommand{\topfraction}{0.9} 
\renewcommand{\textfraction}{0.1}

% Page setup and overhangs
\sloppy
\widowpenalty=10000
\clubpenalty=10000
\hyphenpenalty=75

% Article details for accepted manuscripts will be added by editorial staff
% Omit year if article in press
% Omit number if article under review
\infoteldetails{%
   volume=xx, number=x, month=Month, year=xxxx, firstpage=1, lastpage=10, 
  doi={\href{https://ejournal.ittelkom-pwt.ac.id/index.php/infotel}{10.xxxxx/infotel.vxxix.xxx}},
   received={Month xx, xxxx}, 
   %returned={February 25, 2016},
   revised={Month xx, xxxx},
   accepted={Month xx, xxxx}, }

\newcommand{\mydoi}[1]{\href{http://dx.doi.org/#1}{doi:\protect\detokenize{#1}}}

%\renewcommand{\UrlLeft}{http:\sslash}
%\DeclareUrlCommand\myurl{\def\UrlLeft{}\def\UrlRight{}%
%\urlstyle{tt}}

\urlstyle{rm}
\makeatletter
% Inspired by http://anti.teamidiot.de/nei/2009/09/latex_url_slash_spacingkerning/
% but slightly less kern and shorter underscore
\let\UrlSpecialsOld\UrlSpecials
\def\UrlSpecials{\UrlSpecialsOld\do\/{\Url@slash}\do\_{\Url@underscore}}%
\def\Url@slash{\@ifnextchar/{\kern-.11em\mathchar47\kern-.2em}%
    {\kern-.0em\mathchar47\kern-.08em\penalty\UrlBigBreakPenalty}}
\def\Url@underscore{\nfss@text{\leavevmode \kern.06em\vbox{\hrule\@width.3em}}}
\makeatother

\hypersetup{
colorlinks=true,
linkcolor=black,
citecolor=black,
urlcolor=black
} 

% Add the running author and running title information
\runningauthor{\begin{minipage}{.9\textwidth}\centering Author \textit{et al.}\end{minipage}}
\runningtitle{Running title}

% Document begins
\begin{document}
\setcounter{page}{1}


% Insert your own title
\title{Title of the Article}

% Insert your manuscipts authors, affiliations, and addresses
\author{First Author\textsuperscript{1,*}, Second Author\textsuperscript{2},}
\author{Third Author\textsuperscript{3}}
\affil{\textsuperscript{1,2}Department, Institution, City Zip Code, Country\\
\textsuperscript{3}Department, Institution, City Zip Code, Country}
\affil{*Corresponding email: cauthor@email.ac.id}

\maketitle

% Add 5-10 keywords for every submission
\keywords{keyword 1, keyword 2, keyword 3, keyword 4, keyword 5}

% Add a short abstract of 150-250 words
\begin{abstract}
Abstract max. 250 words.
\end{abstract}

% Your main text begins here. 
\section{Introduction}\label{sec1}
Please, clearly state the introduction, motivation, and contribution of the article here.

\section{Research Method}
This section is not merely for the research method but also can be modified according to the authors' necessity. For example, Related Works, Preliminaries, \textit{etc}.

\section{Results}
\begin{figure}[tbhp]
	\centerline{\includegraphics[scale=0.6]{fig1}}
	\caption{Caption of figure.}
	\label{fig1}
\end{figure}


\section{Discussion}

\begin{table}
	\centering
	\caption{Caption of table}\label{tab1}
	\begin{tabular}{cllll}
		\hline
		\textbf{Head 1} & \textbf{Head 2} & \textbf{Head 3} & \textbf{Head 4} & \textbf{Head 5}\\
		\hline
		Row 1 & Value 1 & Value 2 & Value 3 & Value 4\\
		Row 2 & Value 5 & Value 6 & Value 7 & Value 8\\
	    Row 3 & Value 9 & Value 10 & Value 11 & Value 12\\
		\hline
	\end{tabular}
\end{table}










\section{Conclusion}


\section*{Acknowledgments}
Write the acknowlegment here if any.


%\bibliographystyle{infotelacm}
%\bibliography{infotelexample}


\begin{thebibliography}{00}
%\textit{Basic format for periodical journals:}

\bibitem{b1} R. Fardel, M. Nagel, F. Nuesch, T. Lippert, and A. Wokaun, “Fabrication of organic light emitting diode pixels by laser-assisted forward transfer,” \textit{Appl. Phys. Lett.}, vol. 91, no. 6, Aug. 2007, Art. no. 061103. 

\bibitem{b2} D. Comite and N. Pierdicca, "Decorrelation of the near-specular land scattering in bistatic radar systems," \textit{IEEE Trans. Geosci. Remote Sens.}, early access, doi: 10.1109/TGRS.2021.3072864. (Note: This format is used for articles in early access. The doi must be included.)

\bibitem{b3} H. V. Habi and H. Messer, "Recurrent neural network for rain estimation using commercial microwave links," \textit{IEEE Trans. Geosci. Remote Sens.}, vol. 59, no. 5, pp. 3672-3681, May 2021. [Online]. Available: https://ieeexplore.ieee.org/document/9153027.


\textit{Basic format for conference proceedings examples:}
\bibitem{b4} D. B. Payne and J. R. Stern, “Wavelength-switched passively coupled single-mode optical network,” in \textit{Proc. IOOC-ECOC}, Boston, MA, USA, 1985, pp. 585–590.

\bibitem{b5} D. Ebehard and E. Voges, “Digital single sideband detection for interferometric sensors,” presented at the 2nd Int. Conf. Optical Fiber Sensors, Stuttgart, Germany, Jan. 2-5, 1984.

\bibitem{b6} PROCESS Corporation, Boston, MA, USA. Intranets: Internet technologies deployed behind the firewall for corporate productivity. Presented at INET96 Annual Meeting. [Online]. Available:  http://home.process.com/Intranets/wp2.htp


\textit{Basic format for books examples:}
\bibitem{b7} G. O. Young, “Synthetic structure of industrial plastics,” in \textit{Plastics}, 2nd ed., vol. 3, J. Peters, Ed. New York, NY, USA: McGraw-Hill, 1964, pp. 15–64.

\bibitem{b8} W.-K. Chen, \textit{Linear Networks and Systems}. Belmont, CA, USA: Wadsworth, 1993, pp. 123–135.

\bibitem{b9} Philip B. Kurland and Ralph Lerner, eds., \textit{The Founders’ Constitution}. Chicago, IL, USA: Univ. of Chicago Press, 1987, Accessed on: Feb. 28, 2010, [Online]. Available: http://press-pubs.uchicago.edu/founders/ 

\end{thebibliography}

\end{document}